% Figure: a rational function with a repeated (double) pole, shown as the
% sum of a simple-pole term and a double-pole term. The double-pole term
% dominates near the pole and diverges to the same sign from both sides,
% in contrast with the sign-flipping simple pole of the distinct-pole case.
\begin{figure}[htbp]
\centering
\resizebox{0.82\linewidth}{!}{%
\begin{tikzpicture}
\begin{axis}[
  width=12cm, height=7.4cm,
  xlabel={$x$},
  ylabel={$r(x)$},
  xmin=-1.2, xmax=3.2, ymin=-6, ymax=10,
  xtick={-1,0,1,2,3},
  ytick={-6,-4,-2,0,2,4,6,8,10},
  axis lines=middle,
  tick align=outside,
  every axis plot/.append style={very thick},
  label style={font=\footnotesize},
  tick label style={font=\footnotesize},
  clip=true,
]
  % r(x) = 1 / [ (x-1)^2 (x+1) ]. Double pole at x=1, simple pole at x=-1.
  % Partial fractions: r = A/(x+1) + B/(x-1) + C/(x-1)^2,
  % with A = 1/4, B = -1/4, C = 1/2.
  % Plot in three windows to avoid drawing through the poles.
  \addplot[engblue, domain=-0.85:0.9, samples=200] {1/((x-1)^2*(x+1))};
  \addplot[engblue, domain=1.1:3.2, samples=200] {1/((x-1)^2*(x+1))};
  \addplot[engblue, domain=-1.19:-1.05, samples=120] {1/((x-1)^2*(x+1))};
  % vertical asymptotes
  \addplot[engred, dashed, domain=-6:10, samples=2] coordinates {(1,-6) (1,10)};
  \addplot[engamber, dashed, domain=-6:10, samples=2] coordinates {(-1,-6) (-1,10)};
  \node[font=\footnotesize, text=engred, anchor=south west]
    at (axis cs:1.05, 6.5) {double pole $x=1$};
  \node[font=\footnotesize, text=engamber, anchor=south east]
    at (axis cs:-1.05, 6.5) {simple pole $x=-1$};
  \node[font=\footnotesize, text=engblue, anchor=west]
    at (axis cs:2.1, 1.2) {$r(x)$};
\end{axis}
\end{tikzpicture}%
}%
\caption[A repeated-pole rational function]{The rational function
$r(x) = 1/[(x-1)^2(x+1)]$, which the cover-up-plus-differentiation rule
decomposes into $\tfrac{1/4}{x+1} - \tfrac{1/4}{x-1} +
\tfrac{1/2}{(x-1)^2}$. The double pole at $x=1$ carries the term
$\tfrac{1/2}{(x-1)^2}$, which diverges to $+\infty$ from both sides
because the squared denominator stays positive, so the curve does not
flip sign across the repeated pole. The simple pole at $x=-1$ carries a
single $1/(x+1)$ term and flips sign across the pole in the usual way.
Reading the pole multiplicities off the factored denominator tells the
engineer, before any plotting, which singularities reverse the sign of
the modelled quantity and which do not.}
\label{fig:vol02:ch02:repeated-pole}
\end{figure}
